\chapter{Anatomie d'un portefeuille de stat arb}\label{ch:s1:anatomy-of-a-stat-arb-book}

En moyenne, un livre détient 452 actions en position longue et 495 en position courte. Sa position médiane représente 0.25~\% de son capital, et aucune position n'excède 1~\%.
Sur huit années simulées, son meilleur jour rapporte 1.5~\% et son pire perd 1.25~\%, et son profit est la somme de 1.9 million de jours-positions, dont aucun
ne compte beaucoup isolément. Si l'on retire une seule ligne de sa construction (la neutralité aux facteurs de style), les mêmes signaux donnent un livre dont le meilleur et le pire jour sont
près de quatre fois plus grands, et son P\&L quotidien est, à 99.7~\% de sa variance, un pari sur un seul facteur. Cette partie du livre porte sur les stratégies~; son
premier chapitre porte sur le véhicule dans lequel roulent la plupart des stratégies actions, le \omterm{def:s1:anatomy-of-a-stat-arb-book:statarb}{portefeuille de stat arb} market-neutral~: comment il est construit à partir des éléments du livre~7, ce qui le rend
neutre, comment il est financé, et d'où vient le P\&L de chaque jour. La construction est \texttt{firm.statbook}.

\section{Mille petits paris}

\begin{definition}[Arbitrage statistique, portefeuille de stat arb]\label{def:s1:anatomy-of-a-stat-arb-book:statarb}
L'\emph{arbitrage statistique}\index{arbitrage statistique} (statistical arbitrage) est la négociation systématique de nombreuses petites distorsions relatives de prix entre
titres, chacune peu fiable prise isolément et censée rapporter un peu, de sorte que le rendement du livre provient de leur nombre. Un \emph{portefeuille de stat arb}\index{portefeuille de stat arb} (stat-arb book) est un portefeuille long--short
de centaines ou de milliers d'actions construit de cette façon, neutre au marché et, en général, aux secteurs et aux styles, et reconstruit chaque jour ou plus souvent.
\end{definition}

La loi fondamentale de la gestion active (livre~7, chapitre~15) est le plan d'affaires du \omterm{def:s1:anatomy-of-a-stat-arb-book:statarb}{portefeuille de stat arb}~: un ratio d'information croît avec la racine carrée de la
largeur (breadth), de sorte qu'une compétence trop faible pour être exploitée sur une seule action devient une activité rentable sur mille titres, reconstruite chaque jour. Le prix est tout ce que la loi tient pour acquis~:
les paris doivent être indépendants, ce que les expositions communes d'un livre empêchent, et bon marché, ce que la négociation quotidienne de mille titres n'est pas.

Le livre du chapitre vit sur \texttt{firm.synthmkt} (livre~7, chapitre~5), sur l'ensemble de ses titres cotés, pour les années~3 à~10. Son signal est un mélange à parts égales des deux
prédicteurs que le livre~7 a le mieux mesurés~: le retournement résiduel à un jour et le momentum sur 12--1 mois, tous deux en temps réel, lissés avec une demi-vie de deux jours pour ralentir
ses transactions. À chaque clôture, le signal est rendu neutre par régression sur les expositions du modèle de risque, pondéré pour atteindre une exposition brute de trois fois le capital (1.5
en positions longues, 1.5 en positions courtes), plafonné à 1~\% du capital par titre, et conservé le lendemain à un coût de cinq points de base du poids échangé (le \cref{lst:s1:anatomy-of-a-stat-arb-book:book}).

\section{Neutralité~: dollar, bêta, secteur, facteur}

La neutralité est un ensemble de contraintes sur les expositions du livre (la neutralité au dollar, au bêta et aux facteurs sont les termes du livre~7, chapitre~25). Leur but est que le
P\&L du livre provienne de la partie propre à chaque titre de son signal, et non du marché ni d'un style. Le chapitre construit deux versions du même livre. La première est
neutre au dollar, au bêta du marché et aux dix secteurs industriels, ce que promettent de nombreuses descriptions des livres market-neutral~; la seconde l'est aussi aux
styles taille, valeur et momentum.

\begin{center}
\small
\begin{tabular}{@{}lrr@{}}
\toprule
années 3 à 10, 1~milliard de dollars de capital & neutre au bêta et aux secteurs & également aux styles\\
\midrule
rendement annuel, après coûts et financement & 23.6~\% & 22.5~\%\\
volatilité & 23.6~\% & 4.5~\%\\
ratio de Sharpe & 1.00 & 5.03\\
part de la variance du P\&L issue des facteurs & 99.7~\% & 7.6~\%\\
meilleur et pire jour & $+5.8~\%$, $-4.8~\%$ & $+1.5~\%$, $-1.25~\%$\\
rotation unidirectionnelle quotidienne & 42~\% & 91~\%\\
coûts de transaction annuels & $-10.6~\%$ & $-22.9~\%$\\
\bottomrule
\end{tabular}
\end{center}

Le premier livre est neutre à ce à quoi on lui a demandé d'être neutre, et non à ce qu'il détient. Son signal de momentum lui donne une forte exposition au style momentum,
et la seule contribution du facteur momentum a une volatilité de 23.1~\% par an~: le P\&L spécifique du livre est noyé, et 99.7~\% de sa variance relève du facteur
risque factoriel (la \cref{fig:s1:anatomy-of-a-stat-arb-book:cum}). Neutralisé aussi aux styles, le livre conserve son rendement spécifique et perd le pari factoriel~:
sa volatilité tombe à 4.5~\% et les facteurs expliquent 7.6~\% de sa variance. Le ratio de Sharpe de 5 tient à la simulation, qui introduit un fort retournement et une
dérive persistante (livre~7, chapitre~5)~: un livre réel n'en gagnerait qu'une fraction, comme le montre le bilan public de la section~4. Ce qui se transpose, c'est la forme de la
comparaison. Le livre neutre paie sa neutralité en rotation, 91~\% du capital par jour contre 42~\%, parce que les expositions aux styles bougent chaque jour et doivent être
éliminées par des transactions.

\begin{omfigure}
\begin{tikzpicture}
\begin{axis}[name=a,width=6.8cm,height=5.4cm,xlabel={\small année},ylabel={\small P\&L cumulé (fraction du capital)},tick label style={font=\footnotesize},
  xmin=2,xmax=10,ymin=-1,ymax=4,grid=major,grid style={gray!25},title={\small neutre au bêta et aux secteurs},legend pos=north west,legend style={font=\scriptsize},
  table/col sep=comma]
\addplot[omDef,very thick] table[x=year,y=partial_specific]{figdata/strategies-1/01-anatomy-of-a-stat-arb-book/cum.csv};
\addplot[omThm,thick] table[x=year,y=partial_factor]{figdata/strategies-1/01-anatomy-of-a-stat-arb-book/cum.csv};
\addplot[black!60,thick,dashed] table[x=year,y=partial_total]{figdata/strategies-1/01-anatomy-of-a-stat-arb-book/cum.csv};
\legend{spécifique,facteur,net des coûts et du financement}
\end{axis}
\begin{axis}[at={(a.east)},anchor=west,xshift=1.2cm,width=6.8cm,height=5.4cm,xlabel={\small année},tick label style={font=\footnotesize},
  xmin=2,xmax=10,ymin=-1,ymax=4,grid=major,grid style={gray!25},title={\small également neutre aux styles},table/col sep=comma]
\addplot[omDef,very thick] table[x=year,y=full_specific]{figdata/strategies-1/01-anatomy-of-a-stat-arb-book/cum.csv};
\addplot[omThm,thick] table[x=year,y=full_factor]{figdata/strategies-1/01-anatomy-of-a-stat-arb-book/cum.csv};
\addplot[black!60,thick,dashed] table[x=year,y=full_total]{figdata/strategies-1/01-anatomy-of-a-stat-arb-book/cum.csv};
\end{axis}
\end{tikzpicture}

\omcaption{P\&L quotidien cumulé des deux livres, attribué aux facteurs du modèle de risque et aux rendements spécifiques, et net des coûts de transaction et du financement.
À gauche, les écarts du facteur momentum dominent~; à droite, la courbe factorielle reste proche de zéro. Données~: \texttt{s1\_statbook.run}.}\label{fig:s1:anatomy-of-a-stat-arb-book:cum}
\end{omfigure}

\section{Financement~: trésorerie, levier, emprunt et rebates}

\begin{definition}[Liste des titres difficiles à emprunter]\label{def:s1:anatomy-of-a-stat-arb-book:htb}
La \emph{liste des titres difficiles à emprunter}\index{liste des titres difficiles à emprunter} (hard-to-borrow list) d'un prime broker recense les actions qu'il ne peut prêter aux frais de collatéral général~: les vendre à
découvert coûte des frais d'emprunt plus élevés (livre~1, chapitre~6), peut être refusée et expose au rappel.
\end{definition}

Un livre long--short est un montage de financement autant qu'un portefeuille. Les positions longues sont payées avec le capital du fonds et, au-delà, avec des liquidités empruntées au
prime broker~; les positions courtes sont vendues avec des titres empruntés, et le produit de la vente reste chez le courtier en tant que collatéral, rapportant le taux de rebate (livre~1, chapitre~6),
soit le taux de référence diminué des frais d'emprunt du titre. \texttt{firm.statbook.financing} (le \cref{lst:s1:anatomy-of-a-stat-arb-book:fin}) facture un jour de chacun de ces éléments~: un
solde débiteur de 0.5 du capital (positions longues de 1.5 pour un capital de 1) au taux de référence de 4~\% plus un spread de 0.5~\%, et un produit de vente à découvert de 1.5 au taux de 4~\% diminué des frais,
avec 95~\% des titres à des frais de collatéral général de 0.25~\% et 5~\%, la \omterm{def:s1:anatomy-of-a-stat-arb-book:htb}{liste des titres difficiles à emprunter}, aux frais voisins de 3~\%, non vendus à découvert par le livre. Pour le
livre neutre, le rebate rapporte 5.6~\% du capital par an et le financement de la position longue coûte 2.2~\%~; les frais d'emprunt inclus dans le rebate coûtent 0.4~\%. Le financement net est de $+3.4~\%$,
moins que les 4~\% que rapporterait le capital en trésorerie, et le financement de la position longue et les frais d'emprunt réunis représentent 6.2~\% de l'alpha brut du livre, qui s'élève à 42.0~\%.

Une exposition brute de trois fois le capital n'est pas autorisée par les règles de marge fondées sur la stratégie appliquées aux comptes de particuliers~; c'est le levier d'un livre professionnel, atteint avec
une marge fondée sur le risque (l'\cref{dat:s1:anatomy-of-a-stat-arb-book:margin}).

\begin{dated}{2026-09}{Marge pour un livre d'actions long--short}\label{dat:s1:anatomy-of-a-stat-arb-book:margin}
Aux États-Unis, la Regulation~T (12 CFR 220.12) exige une marge de 50~\% de la valeur d'une position longue sur actions et de 150~\% de la valeur d'une vente à découvert, dont le
produit couvre 100~\%~: les fonds propres d'un livre doivent donc être égaux à la moitié de son brut, de sorte que l'exposition brute est au plus le double du capital~; le livre du chapitre, à trois fois, nécessiterait des fonds propres de
égaux à 1.5 fois son capital. La FINRA Rule~4210 permet aux membres d'utiliser à la place la marge de portefeuille, où les exigences découlent d'un stress de chaque
position (plus ou moins 15~\% pour un titre de capital) plutôt que de pourcentages fixes~; un livre soumis à un stress de 15~\% sur chaque position, sans compensation, a besoin de 0.45
du capital pour un brut de trois. La marge d'entretien fondée sur la stratégie d'une vente à découvert sur une action cotée 5~\$ ou davantage est le plus élevé de 5~\$ par action et de 30~\% de sa valeur.
\end{dated}

\section{Ce que disent les descriptions publiques de tels livres}

Le bilan publié des performances du stat arb est mince et daté, et il dit deux choses. Avellaneda et Lee, qui négocient les résidus contre les composantes principales et
les fonds sectoriels sur le marché américain, ont fait état de ratios de Sharpe après coûts de 1.44 pour 1997--2007, plus élevés avant 2003 et de 0.9 pour 2003--2007, la version ETF
se dégradant de façon semblable après 2002. Khandani et Lo ont documenté l'autre face de l'activité~: beaucoup de livres construits ainsi détenaient des positions semblables, et en août 2007
ils ont perdu ensemble (livre~7, chapitre~28). Le livre de Pedersen sur les stratégies de hedge fund est le récit du praticien à lire à côté d'eux. Les chapitres qui suivent reprennent ses ingrédients un à un~: le retournement
(chapitre~2), les résidus (chapitre~3), les paires (chapitre~4), le momentum (chapitre~5) et le reste de la partie~I.

\section{Le cycle quotidien d'un livre de stat arb}

\begin{definition}[Attribution quotidienne du P\&L]\label{def:s1:anatomy-of-a-stat-arb-book:attribution}
L'\emph{attribution quotidienne du P\&L}\index{attribution quotidienne du P\&L} répartit le P\&L d'un livre pour la journée entre la part expliquée par ses expositions aux facteurs multipliées par les rendements des facteurs du modèle de risque, la part spécifique (ses poids multipliés par les rendements spécifiques des actions),
les coûts de transaction et le financement, de sorte que les parts s'additionnent pour donner le P\&L.
\end{definition}

Le déroulement d'une journée de livre suit un rythme fixe. Après la clôture, les données du jour sont vérifiées et stockées en tenant compte de la date de disponibilité~; le modèle de risque est mis à jour~; les signaux sont calculés,
agrégés et lissés~; le livre est optimisé sous ses contraintes de neutralité et ses limites, et la liste des transactions est transmise à l'exécution pour la séance suivante~; pendant la nuit, le financement
s'accumule~; le lendemain soir, le P\&L de la journée est attribué et rapproché du relevé du prime broker. L'attribution est le moment où un desk voit qu'un livre n'est pas
ce qu'il croit être~: avec les rendements des facteurs $f_t$, les expositions $X$ et les rendements spécifiques $e_t$ du modèle de risque, le rendement de la journée $w^\top r_t = (X^\top w)^\top f_t + w^\top e_t$ tient
exactement, de sorte que la décomposition est une identité, et les contrôles du programme du chapitre confirment que les parts attribuées s'additionnent au P\&L à $10^{-15}$ du capital près.

\begin{omfigure}
\begin{tikzpicture}[st/.style={draw=omDef,fill=omDef!8,rounded corners=2pt,minimum width=1.95cm,align=center,font=\scriptsize,minimum height=1.0cm},
  arr/.style={-{Stealth},thick,omDef}]
\node[st] (d) at (0,0) {clôture~:\\données vérifiées};
\node[st] (r) at (2.35,0) {modèle de risque\\mis à jour};
\node[st] (s) at (4.7,0) {signaux\\agrégés};
\node[st] (o) at (7.05,0) {livre\\optimisé};
\node[st] (e) at (9.4,0) {transactions\\lendemain};
\node[st] (a) at (11.75,0) {P\&L attribué,\\rapproché};
\foreach \x/\y in {d/r,r/s,s/o,o/e,e/a} \draw[arr] (\x) -- (\y);
\draw[omThm,thick,-{Stealth},dashed] (a.south) to[bend left=14] node[below,font=\scriptsize,text=omThm] {le financement s'accumule pendant la nuit~; le cycle suivant recommence à la clôture} (d.south);
\end{tikzpicture}

\omcaption{Le cycle quotidien d'un livre de stat arb.}\label{fig:s1:anatomy-of-a-stat-arb-book:cycle}
\end{omfigure}

\section{Tutoriel~: trois millions de petits paris}\label{tut:s1:anatomy-of-a-stat-arb-book:tut}

\begin{tutorial}
\textbf{Objectif.} Construire les deux livres sur \texttt{firm.synthmkt}, les financer et attribuer leur P\&L quotidien. \textbf{État final~:} le tableau et
la \cref{fig:s1:anatomy-of-a-stat-arb-book:cum}.
\begin{tutsteps}
\item \textbf{Le livre}~: agréger, lisser, neutraliser, plafonner, éviter les ventes à découvert sur les titres difficiles à emprunter.
\omcode{code/strategies-1/01-anatomy-of-a-stat-arb-book/python/s1_statbook.py}{55}{70}{Une clôture de la construction du livre.}\label{lst:s1:anatomy-of-a-stat-arb-book:book}
\item \textbf{Financement et attribution.}
\omcode{code/firm/statbook/firm_statbook.py}{50}{74}{Le financement d'une journée, et l'attribution qui s'additionne.}\label{lst:s1:anatomy-of-a-stat-arb-book:fin}
\item \textbf{Exécuter} \texttt{s1\_statbook.summary} pour les deux livres, \texttt{bets}, \texttt{margin\_example} et \texttt{fig\_statbook.py}.
\end{tutsteps}
\textbf{Pour aller plus loin.} Relever les frais d'emprunt des titres difficiles à emprunter et laisser le livre les vendre à découvert~; neutraliser uniquement le style momentum
et voir quelle part du risque factoriel subsiste.
\end{tutorial}

\section{Construction~: comptabilité du livre de stat arb}\label{bld:s1:anatomy-of-a-stat-arb-book:statbook}

\begin{build}
\bfield{Rôle} La comptabilité de toute stratégie long--short du livre~: financement titre par titre, marge, et une attribution quotidienne qui se rapproche du P\&L.
\bfield{Interface} \texttt{borrow\_fees(n, seed, htb\_share, gc\_fee, htb\_median)}, \texttt{hard\_to\_borrow(fees, threshold)},
\texttt{reg\_t\_equity(long\_mv, short\_mv)}, \texttt{stress\_equity(w, capital, move)}, \texttt{financing(w, capital, rate, long\_spread, fees)}, \texttt{attribution(w, X, f, e, cost, fin)}.
\bfield{Règles} Le financement est imputé titre par titre, chaque jour~; les titres difficiles à emprunter sont signalés avant que l'optimiseur ne les voie~; l'attribution doit
s'additionner au P\&L, sinon la journée n'est pas clôturée.
\bfield{Tests d'acceptation} \texttt{code/firm/statbook/tests/}~: la distribution des frais~; la Regulation~T et l'exigence de stress calculées à la main~; le financement d'une journée calculé à la main~; une
attribution qui s'additionne à $w^\top(Xf + e)$ moins les coûts plus le financement.
\bfield{Pour aller plus loin} Les locates et les rappels jour après jour~; la marge sous un stress de marge de portefeuille avec compensation~; le rapprochement avec un relevé du courtier.
\end{build}

\begin{omsources}
\item M.~Avellaneda et J.-H.~Lee, «~Statistical arbitrage in the US equities market~», \emph{Quantitative Finance} 10(7), 2010.
\item A.~E.~Khandani et A.~W.~Lo, «~What happened to the quants in August 2007~?~», \emph{Journal of Financial Markets} 14(1), 2011.
\item L.~H.~Pedersen, \emph{Efficiently Inefficient}, Princeton University Press, 2015.
\item Regulation~T, 12 CFR 220.12~; FINRA Rule 4210.
\end{omsources}

\section{Exercices}

\begin{exercise}[$\star$]\label{exo:s1:anatomy-of-a-stat-arb-book:1}
Un livre détient 946 positions et négocie chaque jour pendant un an~; si chaque position-jour était un pari indépendant d'écart type 2~\% de la position, quel serait
l'écart type, en unités d'une position, du P\&L cumulé de l'année~?
\end{exercise}

\begin{exercise}[$\star$]\label{exo:s1:anatomy-of-a-stat-arb-book:2}
Le financement net du livre neutre s'élève à 3.42~\% par an~; quel est son coût par rapport au capital laissé en liquidités au taux de l'indice de référence de 4~\%~?
\end{exercise}

\begin{exercise}[$\star$]\label{exo:s1:anatomy-of-a-stat-arb-book:3}
Quels fonds propres la Regulation~T exige-t-elle pour un livre long de 1.5 et short de 1.5 par unité de capital~? Quelle est la plus grande exposition brute qu'elle autorise, et que donne un stress de 15~\%
sans compensation~?
\end{exercise}

\begin{exercise}[$\star\star$]\label{exo:s1:anatomy-of-a-stat-arb-book:4}
Un livre est long de 1.5 et short de 1.5, toutes les ventes à découvert étant soumises aux frais du collatéral général de 0.25~\%, l'indice de référence à 4~\% et le spread long à 0.5~\%. Calculer son financement sur
une année.
\end{exercise}

\begin{exercise}[$\star\star$]\label{exo:s1:anatomy-of-a-stat-arb-book:5}
Pourquoi la neutralisation des styles a-t-elle doublé le turnover du livre~?
\end{exercise}

\begin{exercise}[$\star\star$]\label{exo:s1:anatomy-of-a-stat-arb-book:6}
Montrer que l'attribution $w^\top r = (X^\top w)^\top f + w^\top e$ est une identité lorsque $e$ est le résidu de la régression transversale, et expliquer ce qui la
ferait défaillir.
\end{exercise}

\begin{exercise}[$\star\star\star$]\label{exo:s1:anatomy-of-a-stat-arb-book:7}
\emph{Programmation.} Exécuter \texttt{s1\_statbook.summary} avec une demi-vie de lissage d'un jour, puis de cinq, et rapporter le ratio de Sharpe et les coûts de transaction.
\end{exercise}

\begin{exercise}[$\star\star\star$]\label{exo:s1:anatomy-of-a-stat-arb-book:8}
\emph{Trouver la faille.} «~Notre livre est neutre au marché~: son bêta est nul chaque jour et il est neutre en dollars, donc ses rendements sont de l'alpha pur.~»
\end{exercise}

\section{Problème~: trois millions de petits paris}

\begin{problem}[{Problème du week-end --- un livre, démonté}]\label{pb:s1:anatomy-of-a-stat-arb-book:1}
Les deux livres du chapitre sur \texttt{firm.synthmkt}.

\medskip\noindent\textbf{Partie I --- Le livre.}
\begin{enumerate}
\item Quel est le signal, et comment le transforme-t-on en poids~?
\item Combien de noms le livre neutre détient-il, et quelle est la taille médiane de ses positions~?
\item Pourquoi la largeur compte-t-elle, et que suppose la loi fondamentale que le livre viole~?
\item Combien de positions-jours font ses huit années~?
\end{enumerate}

\medskip\noindent\textbf{Partie II --- Neutralité.}
\begin{enumerate}[resume]
\item Donner les rendements, les volatilités et les ratios de Sharpe des deux livres.
\item Quelle part de la variance de chaque livre provient des facteurs, et de quel facteur dans le premier~?
\item Que coûte au livre neutre sa neutralité~?
\item Pourquoi un ratio de Sharpe de 5 est-il une propriété de la simulation~?
\end{enumerate}

\medskip\noindent\textbf{Partie III --- Financement.}
\begin{enumerate}[resume]
\item Décomposer le financement du livre neutre en taux de rebate, financement de la position longue et frais d'emprunt.
\item Que permet la Regulation~T, et comment un livre atteint-il une exposition brute de trois~?
\item Comment traite-t-on les titres difficiles à emprunter (hard-to-borrow), et combien coûtent-ils~?
\end{enumerate}

\medskip\noindent\textbf{Partie IV --- Le verdict.}
\begin{enumerate}[resume]
\item Énoncer le \emph{résultat nommé}~: la part de la variance du P\&L de l'année expliquée par les expositions aux facteurs dans le livre qui n'était pas neutre aux facteurs de style, et le
poids du financement en fraction de l'alpha brut.
\item Qu'ont trouvé Avellaneda et Lee, et que s'est-il passé après 2002~?
\item Qu'a montré août 2007 sur les livres construits de cette façon~?
\item Décrire le cycle quotidien.
\item Pourquoi faut-il que l'attribution boucle~?
\item Que faudrait-il ajouter pour rendre le livre neutre réaliste~?
\item Comment reconnaître un pari sur facteur déguisé en \omterm{def:s1:anatomy-of-a-stat-arb-book:statarb}{arbitrage statistique} (stat arb)~?
\item Quelle est la première question à poser sur les rendements de tout fonds neutre au marché~?
\item En une phrase~: qu'est-ce qu'un \omterm{def:s1:anatomy-of-a-stat-arb-book:statarb}{portefeuille de stat arb}~?
\end{enumerate}
\end{problem}

\section{Questions d'entretien}

\begin{interviewq}[$\star$ \iqroles{chercheur, trader}]\label{iq:s1:anatomy-of-a-stat-arb-book:1}
Que signifie pour un livre actions long--short d'être neutre au marché, et à quoi reste-t-il exposé~?
\end{interviewq}

\begin{interviewq}[$\star\star$ \iqroles{risque}]\label{iq:s1:anatomy-of-a-stat-arb-book:2}
Comment un livre actions long--short est-il financé, et que peut-il mal tourner dans ce financement~?
\end{interviewq}

\begin{interviewq}[$\star\star$ \iqroles{chercheur}]\label{iq:s1:anatomy-of-a-stat-arb-book:3}
Expliquer la loi fondamentale de la gestion active et pourquoi elle favorise l'\omterm{def:s1:anatomy-of-a-stat-arb-book:statarb}{arbitrage statistique}.
\end{interviewq}

\begin{interviewq}[$\star\star$ \iqroles{chercheur, trader}]\label{iq:s1:anatomy-of-a-stat-arb-book:4}
Les ratios de Sharpe publiés pour le stat arb ont chuté après le début des années 2000. Pourquoi cela pourrait-il être le cas~?
\end{interviewq}

\begin{interviewq}[$\star\star$ \iqroles{développeur, risque}]\label{iq:s1:anatomy-of-a-stat-arb-book:5}
Comment concevriez-vous l'\omterm{def:s1:anatomy-of-a-stat-arb-book:attribution}{attribution quotidienne du P\&L} d'un livre de mille noms~? Qu'est-ce qui doit se réconcilier~?
\end{interviewq}

\begin{interviewq}[$\star\star\star$ \iqroles{chercheur}]\label{iq:s1:anatomy-of-a-stat-arb-book:6}
Un livre a une exposition brute $G$ avec $N$ positions égales, une volatilité spécifique $\sigma$ par titre et des rendements spécifiques indépendants. Dériver sa volatilité spécifique et le
ratio d'information pour un edge $\alpha$ par titre~; comment évolue-t-il avec $N$~?
\end{interviewq}
